% Технічне завдання (ГОСТ 19.201-78), за шаблоном кафедри ч1 2026.
\documentclass{kpi-thesis}
\usepackage{kpi-titles}
% Дані роботи — єдине джерело для всіх документів (титулка, завдання, шифри, штампи).
% Тема — дослівно як у наказі.
\kpiset{
  title             = {Тема дипломного проєкту},
  group             = {ІП-XX},
  number            = {NN},          % номер у списку групи: шифр КПІ.ІП-XXNN...
  devcode           = {045440},      % вид розробки, методичка §5.4 (045440 — вебзастосунки)
  gender            = {m},           % f — «Виконала», «студентка»
  student           = {Прізвище Ім'я По батькові},
  student-dative    = {Прізвищу Імені По батькові},
  student-short     = {Ім'я ПРІЗВИЩЕ},
  student-stamp     = {Прізвище І.П.},
  advisor           = {посада, науковий ступінь, вчене звання, Прізвище І. П.},
  advisor-full      = {Прізвище Ім'я По батькові, науковий ступінь, вчене звання},
  advisor-short     = {Ім'я ПРІЗВИЩЕ},
  advisor-stamp     = {Прізвище І.П.},
  consultant        = {посада, Прізвище Ім'я По батькові},
  normcontrol-short = {Ім'я ПРІЗВИЩЕ},
  normcontrol-stamp = {Прізвище І.П.},
  reviewer          = {посада, науковий ступінь, вчене звання, Прізвище І. П.},
  year              = {2027},
  order-date        = {\todo{«3» червня 2027 р.}},
  order-number      = {\todo{1234-с}},
  deadline          = {\todo{«14» червня 2027 року}},
  issued            = {\todo{«15» березня 2027 року}},
  head-short        = {Едуард~ЖАРІКОВ},    % завідувач кафедри — звірте з поточним
  head-stamp        = {Жаріков~Е.В.},
  % Типові значення класу; змінити, якщо інші:
  % specialty       = {121~Інженерія~програмного~забезпечення},
  % program         = {Інженерія~програмного~забезпечення~інформаційних~систем},
}
% Терміни календарного плану (завдання, сторінка 2).
\renewcommand\kpiCalendarRows{%
  1 & Вивчення рекомендованої літератури & & \\ \hline
  2 & Аналіз існуючих методів розв'язання задачі & & \\ \hline
  3 & Постановка та формалізація задачі & & \\ \hline
  4 & Розробка інформаційного забезпечення & & \\ \hline
  5 & Алгоритмізація задачі & & \\ \hline
  6 & Обґрунтування вибору використаних технічних засобів & & \\ \hline
  7 & Розробка програмного забезпечення & & \\ \hline
  8 & Налагодження програми & & \\ \hline
  9 & Виконання графічних документів & & \\ \hline
  10 & Оформлення пояснювальної записки & & \\ \hline
  11 & Подання ДП на попередній захист & & \\ \hline
  12 & Подання ДП рецензенту & & \\ \hline
  13 & Подання ДП на основний захист & & \\ \hline
}


\begin{document}
\kpiDocTitle{specification}
\mainmatter
\tableofcontents

\section{Найменування та галузь застосування}
Назва розробки: \kpiTitle.

Галузь застосування: \todo{на розробку якого ПЗ поширюється ТЗ і для чого воно використовується}.

\section{Підстава для розробки}
Підставою для розробки є завдання на дипломне проєктування, затверджене кафедрою інформатики та програмної інженерії Національного технічного університету України «Київський політехнічний інститут імені Ігоря Сікорського».

\section{Призначення розробки}
Розробка призначена для \todo{функціональне та експлуатаційне призначення}.

Метою розробки є \todo{показники якості або процес досягнення якості}.

\section{Вимоги до програмного забезпечення}
\subsection{Вимоги до функціональних характеристик}
Програмне забезпечення повинно забезпечувати виконання наступних основних функцій.

\subsubsection{Користувацького інтерфейсу}
\todo{Обов'язково прототипи екранних форм (UX). 5 і більше функціональних задач верхнього рівня з вимогами за ролями.}

\subsubsection{Для користувача}
\begin{itemize}
\item \todo{функція};
\end{itemize}

\subsubsection{Для адміністратора системи}
\begin{itemize}
\item \todo{функція};
\end{itemize}

\subsubsection{Додаткові вимоги}
\todo{...}

\subsection{Вимоги до надійності}
\todo{...}

\subsection{Умови експлуатації}
\subsubsection{Вид обслуговування}
\todo{...}
\subsubsection{Обслуговуючий персонал}
\todo{...}

\subsection{Вимоги до складу і параметрів технічних засобів}
\todo{...}

\subsection{Вимоги до інформаційної та програмної сумісності}
\subsubsection{Вимоги до вхідних даних}
\todo{...}
\subsubsection{Вимоги до вихідних даних}
\todo{...}
\subsubsection{Вимоги до мови розробки}
\todo{...}
\subsubsection{Вимоги до середовища розробки}
\todo{...}
\subsubsection{Вимоги до представлення вихідних кодів}
\todo{...}

\subsection{Вимоги до маркування та пакування}
Вимоги до маркування та пакування не висуваються.

\subsection{Вимоги до транспортування та зберігання}
Вимоги до транспортування та зберігання не висуваються.

\subsection{Спеціальні вимоги}
\todo{...}

\section{Вимоги до програмної документації}
\subsection{Попередній склад програмної документації}
У склад супровідної документації повинні входити такі документи:
\begin{itemize}
\item пояснювальна записка;
\item технічне завдання;
\item текст програми;
\item програма та методика тестування;
\item керівництво користувача;
\item графічний матеріал.
\end{itemize}

\subsection{Спеціальні вимоги до програмної документації}
\todo{...}

\section{Стадії і етапи розробки}
\todo{Етапи з термінами, узгоджені з календарним планом.}

\section{Порядок контролю та приймання}
\todo{Склад і обсяг приймально-здавальних випробувань.}

\end{document}
